Reduction is a powerful paradigm in both algorithms and complexity.
In algorithms, a dual view of a problem can sometimes provide a window
to a more efficient algorithm. With this as the main motivation, we
discuss the problem of convex optimization given different types of
access to the instance at hand. Despite appearing to be rather different,
all these perspectives will turn out to be equivalent, i.e., reducible
to each other, with only polynomial overhead. The main tool is duality.
This chapter elaborates on these ideas. For the classical treatment
of oracle models, the reader is referred to \cite{grotschel2012geometric}.
For the most efficient known reductions, which are conjectured to
be optimal up to logarithmic factors, see \cite{lee2020efficient}. 

\section{Equivalences between Oracles}

\subsection{Oracles for Convex Sets}

Gr�tschel, Lov\'asz and Schrijver \cite{grotschel2012geometric}
defined five different oracles to access convex sets, showed they
are equivalent and used them to get polynomial-time algorithms for
a variety of combinatorial problems (including the first ones in many
cases). 

Here are four basic oracles\footnote{We also omit the fifth oracle VIOL defined by \cite{grotschel2012geometric},
which checks whether the convex set satisfies a given inequality or
gives a violating point in the convex set, since this is equivalent
to the OPT oracle below up to a logarithmic factor.} for a convex set $K\subseteq\R^{n}$. Later we will allow for error
parameters in each oracle, i.e., approximate versions of all the oracles
below. We begin with exact oracles for simplicity.
\begin{defn}[Membership Oracle (MEM)]
\label{def:MEM}Queried with a vector $y\in\Rn$, the oracle either

\begin{itemize}
\item asserts that $y\in K$, or
\item asserts that $y\notin K$.
\end{itemize}
\end{defn}

\begin{defn}[Separation Oracle (SEP)]
\label{def:SEP}Queried with a vector $y\in\Rn$, the oracle either

\begin{itemize}
\item asserts that $y\in K$, or
\item finds a unit vector $c\in\Rn$ such that $c^{T}x\leq c^{T}y$ for
all $x\in K$.
\end{itemize}
\end{defn}

\begin{defn}[Validity Oracle (VAL)]
\label{def:VAL}Queried with a unit vector $c\in\Rn$, the oracle
either\footnote{We use a slightly different definition than \cite{grotschel2012geometric}
for clarity. }

\begin{itemize}
\item outputs $\max_{x\in K}c^{\top}x$, or
\item asserts that $K$ is empty.
\end{itemize}
\end{defn}

\begin{defn}[Optimization Oracle (OPT)]
\label{def:OPT}Queried with a unit vector $c\in\Rn$, the oracle
either

\begin{itemize}
\item finds a vector $y\in K$ and $c^{T}x\leq c^{T}y$ for all $x\in K$,
or
\item asserts that $K$ is empty.
\end{itemize}
\end{defn}

According to the definitions, the separation oracle gives more information
than the membership oracle and the optimization oracle gives more
information than the validity oracle. Depending on the problem, usually
one of the oracles will be the preferred or more natural way to access
the convex set. For example, for the polytope given by $\{Ax\geq b\}$,
the separation oracle is the preferred way because the membership
oracle takes as much time as the separation oracle, and both validity
and optimization involve solving a linear program. On the contrary,
the preferred oracle for the convex set $\conv(\{a_{i}\})$ is the
optimization oracle because $\max_{x\in\conv(\{a_{i}\})}\theta^{\top}x$
can be solved by only checking $x=a_{i}$ for each $i$. In combinatorial
optimization, many polytopes have exponentially many vertices and
constraints but one can use combinatorial structure to solve the optimization
problem efficiently.
\begin{example}[Zonotope]
 Imagine you have $n$ machines. Each machine uses resources $(r_{1},r_{2},\cdots,r_{k})\in\R^{k}_{+}$
and produces goods $(g_{1},g_{2},\cdots,g_{\ell})\in\R^{\ell}_{+}$.
We represent a machine simply by the vector $a=(r_{1},r_{2},\cdots,r_{k},g_{1},g_{2},\cdots,g_{\ell})$.
Say we have vectors $a_{1},a_{2},\cdots,a_{n}\in\R^{k+\ell}$, then
the polytope
\[
K=\left\{ \sum^{n}_{i=1}\lambda_{i}a_{i}\in\R^{k+\ell}|0\leq\lambda_{i}\leq1\text{ for all }i\right\} 
\]
represents the set of feasible (input, output) pairs. Many production
problems can then be written as
\[
\min_{z\in K}f(z).
\]
\end{example}

\begin{xca}
Evaluate the complexity of each of the basic oracles for the Zonotope
problem above.
\end{xca}

\begin{example}
The spanning tree polytope of an undirected graph $G=(V,E)$ is given
by
\[
P=\{x\in\R^{\left|E\right|}_{+}|\sum_{e\in E}x_{e}=|V|-1,\sum_{(u,v)\in E\cap(S\times S)}x_{(u,v)}\leq|S|-1\quad\forall S\subseteq V,\left|S\right|\ge1\}.
\]
The extreme points of this polytope can be shown to exactly correspond
to the indicator vectors of spanning trees of $G$. Thus, the optimization
oracle in this case is to simply find a maximum cost spanning tree.
\end{example}

\begin{xca}
Design a membership oracle for the spanning tree polytope.
\end{xca}

\begin{figure}[h]
\centering
\begin{tikzpicture}[node distance = 1.8cm]
\tikzstyle{block} = [rectangle, draw, text width=6.5em, text centered, rounded corners, minimum height=4em]
\tikzstyle{line} = [draw, -tonew,arrowhead=0.085cm]
\tikzstyle{dot_line} = [draw, -tonew,dashed,arrowhead=0.085cm]
% Place nodes
   \node [block] (opt) {\scriptsize $OPT(K) = \partial {\delta_K^*}$};
   \node [block, below of=opt, node distance=1.8cm] (val) {\scriptsize $VAL(K) = {\delta_K^*}$};
   \node [block, right of=opt, node distance=3cm] (sepp) {\scriptsize $SEP(K) = \partial{\delta_K}$};
   \node [block, below of=sepp, node distance=1.8cm] (mem) {\scriptsize $MEM(K) = {\delta_K}$};
   \node [left of=opt, node distance=3cm] (viol) {};
   \node [below of=viol, node distance=1.5cm] {$\xrightarrow{\quad\quad} \mbox{ is } \tilde{O}(1)$};
   \node [below of=viol, node distance=2.1cm] {$\xdashrightarrow{\quad\quad} \mbox{ is } \tilde{O}(n)$}; 
   % Draw edges
%   \path [line] (opt) -- (viol);
%   \path [line] (viol) -- (opt);
   \path [line,transform canvas={xshift=-0.25cm}] (sepp) -- (mem);
   \path [line,transform canvas={xshift=-0.25cm}] (opt) -- (val);
   \path [dot_line,transform canvas={xshift=0.25cm}] (mem) -- (sepp);
   \path [dot_line,transform canvas={xshift=0.25cm}] (val) -- (opt);
   \path [dot_line,transform canvas={yshift=-0.25cm}] (opt) -- (sepp);
   \path [dot_line,transform canvas={yshift=0.25cm}] (sepp) -- (opt);  
\end{tikzpicture}

\caption{The relationships among the four oracles for convex sets. The arrows
are implications. \label{fig:oracles}}
\end{figure}


\subsection{Duals of Convex Sets}

\begin{figure}[h]
\centering{}\includegraphics[width=0.75\columnwidth]{Figures/Section_4\lyxdot 1\lyxdot 2}\caption{Insights into the Dual. (a) K contains the origin. (b) K does not
contain the origin.}
\end{figure}

The dual or \emph{polar} of a convex set $K$ is defined as the following
convex set:
\[
K^{*}=\left\{ y:\forall x\in K,\left\langle x,y\right\rangle \le1\right\} .
\]
One might imagine that the polar operator is an involution, i.e.,
the dual of the dual is the original set. Although this is not always
true, the set $(K^{*})^{*}$ still has a simple description and under
mild conditions we have $(K^{*})^{*}=K$.
\begin{prop}
\label{prop:K**} Let $K\subseteq$$\,\R^{n}$ be a convex set. Then
\[
(K^{*})^{*}=\mathrm{cl(conv(\mathit{K}\cup\left\{ 0\right\} ))}.
\]
\end{prop}

\begin{proof}
We first prove that $K^{**}\supseteq\mathrm{cl(conv(\mathit{K}\cup\left\{ 0\right\} ))}$.
The polar of any convex set is closed and convex since it is the intersection
of halfspaces, which are closed and convex. Thus, to prove this part,
it is sufficient to prove that $K\subseteq K^{**}$ and $0\in K^{**}$.
Let $x\in K$. Then, for every $y\in K^{*}$, it holds that $\langle x,y\rangle\leq1$,
by the definition of the polar. Note that these inequalities are precisely
the constraints of $K^{**}$, thus $x\in K^{**}$. Moreover, $\langle0,y\rangle=0\leq1$,
for every $y\in K^{*}$, and so $0\in K^{**}$.

We now prove the other direction. Suppose there is a vector $y\in K^{**}\setminus\mathrm{cl(conv(\mathit{K}\cup\left\{ 0\right\} ))}$.
By Theorem \ref{thm:sep}, there is $\theta$ such that
\[
\theta^{\top}y>\sup_{x\in cl(conv(\mathit{K}\cup\left\{ 0\right\} ))}\theta^{\top}x.
\]
Since $0$ is feasible, the RHS is at least $0$. Thus, $\theta^{\top}y>0$.
Scale $\theta$ so that 
\[
\theta^{\top}y>1\geq\sup_{x\in cl(conv(\mathit{K}\cup\left\{ 0\right\} ))}\theta^{\top}x.
\]
Hence, $\theta\in K^{*}$ since $\theta^{\top}x\leq1$ for every $x\in K$.
However, since $\theta^{\top}y>1$, we have $y\not\in K^{**}$, a
contradiction.
\end{proof}
\begin{cor}
Let $K\subseteq$$\,\R^{n}$ be a closed convex set. Then $(K^{*})^{*}=K$
if and only if $0\in K$.
\end{cor}


\subsection{Oracles for Convex Functions}

Now, we generalize the oracles for convex sets to convex functions.
The membership oracle and separation oracle can be generalized as
follows
\begin{defn}[Evaluation Oracle (EVAL)]
\label{def:EVAL}Queried with a vector $y$, the oracle outputs $f(y)$.
\end{defn}

The next oracle generalizes gradients to subgradients so that they
can be computed for general convex functions. Any vector output by
the oracle is a subgradient of the function at the query point.
\begin{defn}[Subgradient Oracle (GRAD)]
\label{def:GRAD}Queried with a vector $y$, the oracle outputs $f(y)$
and a vector $g\in\Rn$ such that
\begin{equation}
f(x)\geq f(y)+g^{\top}(x-y)\text{ for all }x\in\Rn\label{eq:grad_guarantee}
\end{equation}
or outputs ``$f(y)$ undefined.''
\end{defn}

To generalize the validity oracle, we define the convex (Fenchel)
conjugate of $f$:
\begin{defn}[Convex Conjugate]
For any function $f$, we define the convex conjugate 
\[
f^{*}(\theta)\defeq\sup_{x\in\R^{n}}\theta^{\top}x-f(x).
\]
\end{defn}

Note that $f^{*}$ is convex because it is the supremum of linear
functions. Also, we have $f^{*}(0)=-\inf_{x\in\R^{n}}f(x)$. Note
that\footnote{Recall that $\delta_{C}(x)=0$ if $x\in C$ and $+\infty$ otherwise.}
$\delta^{*}_{K}(c)=\sup_{x\in K}c^{\top}x$. Therefore, the validity
oracle for $\delta_{K}$ is simply the evaluation oracle for $\delta^{*}_{K}$.
The following lemma shows that the optimization oracle is simply the
(sub)gradient oracle for $\delta^{*}_{K}$. We use $\nabla f$ to
represent subgradient. 
\begin{lem}
\label{lem:gradient_conjugate}For any continuous function $f$ with
$f^{*}$ differentiable at $\theta$, we have that $\nabla f^{*}(\theta)=\arg\max_{x}\theta^{\top}x-f(x)$.
\end{lem}

\begin{proof}
First we observe that the supremum is achieved. Fix $\theta$. Let
$g_{\theta}(x)=\theta^{\top}x-f(x)$. We assume $\sup_{x}g_{\theta}(x)=f^{*}(\theta)$
is finite. Let $\epsilon>0$. Then, $g_{\theta}(x)$ is a continuous
function, so $S=g^{-1}_{\theta}([f^{*}(\theta)-\epsilon,\infty))$
is a closed set. The set $S$ is not empty because there is some $x$
for which $g_{\theta}(x)\ge\sup_{z}g_{\theta}(z)-\epsilon$. Now suppose
for a contradiction that $S$ is not bounded. Then, there exists a
sequence $x_{i}\in\R^{n}$ such that $\norm{x_{i}}\ge i$ and $g_{\theta}(x_{i})\ge f^{*}(\theta)-\epsilon$
for all $i$. By the compactness of the unit sphere, we may assume
by taking a subsequence that $x_{i}/\lVert x_{i}\rVert\to u$ for
some unit vector $u\in\R^{n}$. Then, $u^{\top}x_{i}/\lVert x_{i}\rVert\to u^{\top}u=1$,
and since $\lVert x_{i}\rVert\to\infty$, we have $u^{\top}x_{i}\to\infty$.
Choose $t>0$ small enough that $f^{*}(\theta+tu)<\infty$. Since
$g_{\theta+tu}(x_{i})=g_{\theta}(x_{i})+tu^{\top}x_{i}\ge f^{*}(\theta)-\epsilon+tu^{\top}x_{i}$,
this means that $g_{\theta+tu}(x_{i})\to\infty$, contradicting $f^{*}(\theta+tu)$
being finite. Hence, $S$ is bounded and thus compact, so $g_{\theta}(x)$
attains its maximum in $S$, and thus in $\R^{n}$.

Let $x_{\theta}\in\arg\sup_{x}\theta^{\top}x-f(x)$. By definition,
we have that $f^{*}(\theta)=\theta^{\top}x_{\theta}-f(x_{\theta})$
and that $f^{*}(\eta)\geq\eta^{\top}x_{\theta}-f(x_{\theta})$ for
all $\eta$. Therefore, 
\[
f^{*}(\eta)\geq f^{*}(\theta)+x^{\top}_{\theta}(\eta-\theta)\text{ for all }\eta.
\]
Therefore, $x_{\theta}\in\nabla f^{*}(\theta)$.
\end{proof}
Note that this lemma shows that $\nabla\delta^{*}_{K}(\theta)=\arg\max_{x\in K}\theta^{\top}x$.
So, the gradient oracle for $\delta^{*}_{K}$ is exactly the optimization
oracle for $K$.

Figure \ref{fig:rel} shows the current best reduction between these
four function oracles. Note that according to our definition of the
GRAD oracle, it also outputs the function value (EVAL). It is not
hard to show that you can use $n+1$ calls to evaluation oracle to
compute one gradient approximately (using finite difference) and that
you can use $\tilde{O}(n)$ calls to gradient oracle of $f$ to compute
one gradient of $f^{*}$ (using the cutting plane method). In the
next section (Sec. \ref{sec:From-EVAL-to}), we will see the remaining
reductions. However, it is still an open question if all of these
reductions are the best possible. 
\begin{problem*}
Prove that it takes $\Omega(n^{2})$ calls to the evaluation oracle
of $f$ to compute the gradient of $f^{*}$.
\end{problem*}
\begin{figure}[h]
\centering
\begin{tikzpicture}[node distance = 1.8cm]
\tikzstyle{block} = [rectangle, draw, text width=6.5em, text centered, rounded corners, minimum height=4em]
\tikzstyle{line} = [draw, -tonew,arrowhead=0.085cm]
\tikzstyle{dot_line} = [draw, -tonew,dashed,arrowhead=0.085cm]
% Place nodes
   \node [block] (opt) {\scriptsize $GRAD(f^*)$};
   \node [block, below of=opt, node distance=1.8cm] (val) {\scriptsize $EVAL(f^*)$};
   \node [block, right of=opt, node distance=3cm] (sepp) {\scriptsize $GRAD(f)$};
   \node [block, below of=sepp, node distance=1.8cm] (mem) {\scriptsize $EVAL(f)$};
   \node [left of=opt, node distance=3cm] (viol) {};
   \node [below of=viol, node distance=1.5cm] {$\xrightarrow{\quad\quad} \mbox{ is } \tilde{O}(1)$};
   \node [below of=viol, node distance=2.1cm] {$\xdashrightarrow{\quad\quad} \mbox{ is } \tilde{O}(n)$}; 
   %\node [below of=viol, node distance=1.5cm] {$\tilde{O}(1) \xrightarrow{\quad\quad}$};
   %\node [below of=viol, node distance=2.1cm] {$\tilde{O}(n) \xdashrightarrow{\quad\quad}$}; 
   % Draw edges
   \path [line,transform canvas={xshift=-0.25cm}] (sepp) -- (mem);
   \path [line,transform canvas={xshift=-0.25cm}] (opt) -- (val);
   \path [dot_line,transform canvas={xshift=0.25cm}] (mem) -- (sepp);
   \path [dot_line,transform canvas={xshift=0.25cm}] (val) -- (opt);
   \path [dot_line,transform canvas={yshift=-0.25cm}] (opt) -- (sepp);
   \path [dot_line,transform canvas={yshift=0.25cm}] (sepp) -- (opt);  
\end{tikzpicture}

\caption{This illustrates the relationships of oracles for a convex function
$f$ and its convex conjugate $f^{*}$. \label{fig:rel}}
\end{figure}


\subsection{Convex Conjugate and Equivalences}

Here are some examples of conjugates.
\begin{xca}
Show that
\end{xca}

\begin{itemize}
\item The conjugate of $f(x)=\frac{1}{p}\|x\|^{p}_{p}$ is $f^{*}(\theta)=\frac{1}{q}\|\theta\|^{q}_{q}$
where $\frac{1}{p}+\frac{1}{q}=1$ (with $1<p,q<\infty)$.
\item The conjugate of $f(x)=\left\langle a,x\right\rangle -b$ is $f^{*}(\theta)=\begin{cases}
b, & \theta=a\\
+\infty & \mbox{otherwise}
\end{cases}$.
\item The conjugate of $f(x)=\sum_{i}e^{x_{i}}$ is $f^{*}(\theta)=\begin{cases}
\sum_{i}\theta_{i}\log\theta_{i}-\theta_{i} & \text{if }\theta_{i}\geq0\text{ for all }i\\
+\infty & \mbox{otherwise}
\end{cases}.$
\end{itemize}
%
\begin{xca}
For any $x,\theta$, we have $\theta^{\top}x\leq f(x)+f^{*}(\theta)$.
\end{xca}

\begin{xca}
Prove that the gradient of $f$ is $L$-Lipschitz if and only if $f^{*}$
is $\frac{1}{L}$ strongly convex.
\end{xca}

The following lemma shows that one can recover $f$ from $f^{*}$.
\begin{lem}[Involution property]
\label{lem:involution}For any convex function $f$ with a closed
epigraph, we have that $f^{**}=f$.
\end{lem}

\begin{proof}
Since $\epi f$ is a closed, convex set, Corollary \ref{cor:convex_function_intersection}
shows that it must be an intersection of halfspaces, i.e., a set ${\cal H}$
such that 
\[
f(x)=\sup_{(\theta,b)\in\mathcal{H}}\theta^{\top}x-b
\]
where $\mathcal{H}$ is the set of supporting planes of $\epi f$
and contains all affine lower bounds on $f$, namely, $f(x)\geq\theta^{\top}x-b$
for all $x$. Alternatively, we can write
\[
\mathcal{H}=\left\{ (\theta,b):\,\forall x,\,f(x)\ge\theta^{\top}x-b\right\} ,\qquad\epi(f)=\bigcap_{(\theta,b)\in\mathcal{H}}\left\{ (x,t):t\ge\theta^{\top}x-b\right\} .
\]
For a fixed $\theta$, any feasible $b$ satisfies $b\geq\theta^{\top}x-f(x)$
for all $x$. So, the smallest feasible value satisfies 
\[
b^{*}=\sup_{x}\theta^{\top}x-f(x)=f^{*}(\theta).
\]
Hence, 
\[
f(x)=\sup_{(\theta,b)\in\mathcal{H}}\theta^{\top}x-b^{*}=\sup_{\theta}\theta^{\top}x-f^{*}(\theta)=f^{**}(x).
\]
\end{proof}
\begin{xca}
Let $f$ be a convex function with closed epigraph. Show that $f=f^{*}$
if and only if $f(x)=\frac{\norm x^{2}}{2}$.
\end{xca}

Since we can use the gradient, $\nabla f$, to compute the gradient
of the dual, $\nabla f^{*}$ (via cutting plane method), the involution
property shows that we can do the reverse -- use $\nabla f^{*}$
to compute $\nabla f$. Going back to the example about $\conv(\{a_{i}\})$,
since we know how to compute $\max_{x\in\conv(\{a_{i}\})}\theta^{\top}x=\delta^{*}_{\conv(\{a_{i}\})}(\theta)$,
this reduction gives us a way to separate $\conv(\{a_{i}\})$, or
equivalently, to compute the (sub)gradient of $\delta^{*}_{\conv(\{a_{i}\})}$.
This is formalized in the next exercise.
\begin{xca}
Show how to implement the separation oracle SEP for a convex set $K$
given access to an optimization oracle OPT for $K$.
\end{xca}

Recall that for any linear space $X$, $X^{*}$ denotes the dual space,
i.e., the set of all linear functions on $X$ and that under mild
assumptions\footnote{The dual of the dual of a vector space $X$ is isomorphic to $X$
if and only if $X$ is a finite-dimensional vector space.}, we have $X^{**}=X$. Therefore, there are two natural ``coordinate
systems'' to record a convex function, the primal space $X$ and
the dual space $X^{*}$. Under these ``coordinate systems'', we
have the dual functions $f$ and $f^{*}$.
\begin{xca}[Order reversal]
 Show that $f\geq g\iff f^{*}\leq g^{*}$ (both for all $x\in\R^{n}$).
\end{xca}

\begin{rem*}
Interestingly, the convex conjugate is the \emph{unique} transformation
on convex functions that satisfies both involution and order reversal. 
\end{rem*}
\begin{thm}[\cite{artstein2009concept}]
Given a transformation $\mathcal{T}$ that maps the set of lower
semi-continuous\footnote{See Def.\ref{def:lower-semi}.} convex functions
onto itself such that $\mathcal{T}\mathcal{T}\phi=\phi$ and $\phi\leq\psi\implies\mathcal{T}\phi\geq\mathcal{T}\psi$
for all lower-semi-continuous convex functions $\phi$ and $\psi$.
Then, $\mathcal{T}$ is essentially the convex conjugate, namely,
there is an invertible symmetric linear transformation $B$, a vector
$v_{0}$ and a constant $C_{0}$ such that 
\[
(\mathcal{T}\phi)(x)=\phi^{*}(Bx+v_{0})+v^{\top}_{0}x+C_{0}.
\]
\end{thm}

In combinatorial optimization, many convex sets are given by the convex
hull for some discrete objects. In many cases, the only known way
to do the separation is via such reductions. In this chapter, we will
study the following general theorem showing that optimization can
be reduced to membership/evaluation with a quadratic overhead in dimension
for the number of oracle queries. 
\begin{thm}
Let $K$ be a convex set specified by a membership oracle, a point
$x_{0}\in\R^{n}$, and numbers $0<r<R$ such that $B(x_{0},r)\subseteq K\subseteq B(x_{0},R)$.
For any convex function $f$ given by an evaluation oracle and any
$\epsilon>0$, there is a randomized algorithm that computes a point
$z\in B(K,\epsilon)$ such that 
\[
f(z)\le\min_{x\in K}f(x)+\epsilon\left(\max_{x\in K}f(x)-\min_{x\in K}f(x)\right)
\]
with constant probability using $O\left(n^{2}\log^{2}\left(\frac{nR}{\epsilon r}\right)\right)$
calls to the membership oracle and evaluation oracle and $O(n^{3}\log^{O(1)}\left(\frac{nR}{\epsilon r}\right))$
total arithmetic operations. 
\end{thm}


\section{Gradient from Evaluation via Finite Difference\label{sec:From-EVAL-to}}

When the function $f$ is in $\mathcal{C}^{1}$, we can approximate
the gradient of a function by calculating a finite difference for
sufficiently small $h$:
\[
\frac{\partial f}{\partial x_{i}}=\frac{f(x+he_{i})-f(x)}{h}+O(h)
\]
which only takes $n+1$ calls to the evaluation oracle (for computing
$f(x),f(x+he_{1}),\cdots,f(x+he_{n})$). The only issue is that the
convex function may not be differentiable. However, any convex Lipschitz
function is twice differentiable almost everywhere (see the proof
below). Therefore, we can simply perturb $x$ with random noise, then
apply a finite difference. To see the idea more precisely, we first
observe that the norm of the Hessian can be bounded in expectation
for a Lipschitz function. Note that this is Lipschitzness of the function,
not its gradient. The proof below uses the basic fact that the gradient
is defined almost everywhere for Lipschitz functions. 
\begin{lem}
For any $L$-Lipschitz convex function $f$ defined in a unit ball,
we have $\nabla^{2}f(x)$ exists almost everywhere and that $\E_{x\in B(0,1)}\|\nabla^{2}f(x)\|_{F}\leq nL$.
\end{lem}

\begin{proof}
The existence almost everywhere is classical (Alexandrov Theorem,
see e.g., \cite{constantin2018convex}; see also Rademacher's theorem
about the existence of the derivative almost everywhere for Lipschitz
functions). We will only prove the part $\E_{x\in B(0,1)}\|\nabla^{2}f(x)\|_{F}\leq nL$.
Since $\nabla^{2}f\succeq0$ (where defined), we have $\|\nabla^{2}f(x)\|_{F}\leq\tr\nabla^{2}f(x)$.
Therefore, 
\[
\int_{B(0,1)}\|\nabla^{2}f(x)\|_{F}dx\leq\int_{B(0,1)}\tr\nabla^{2}f(x)dx=\int_{B(0,1)}\Delta f(x)dx.
\]
Using Stokes' Theorem, and letting $\nu(x)$ be the normal vector
at $x$, we have 
\[
\int_{B(0,1)}\Delta f(x)dx=\int_{\partial B(0,1)}\left\langle \nabla f(x),\nu(x)\right\rangle dx\leq|\partial B(0,1)|\cdot L.
\]
Hence, we have
\[
\E_{x\in B(0,1)}\|\nabla^{2}f(x)\|_{F}\leq\frac{|\partial B(0,1)|}{|B(0,1)|}L=nL.
\]
\end{proof}
To turn this into an algorithm, we need to develop it a bit further.
\begin{lem}[\cite{lee2017efficient}]
\label{lem:convex_almost_flat}Let $B_{\infty}(x,r)=\{y:\ \norm{x-y}_{\infty}\leq r\}$.
For any $0<r_{2}\leq r_{1}$ and any convex function $f$ defined
on $B_{\infty}(x,r_{1}+r_{2})$ with $\norm{\nabla f(z)}_{\infty}\leq L$
for any $z\in B_{\infty}(x,r_{1}+r_{2})$ we have 
\[
\E_{y\in B_{\infty}(x,r_{1})}\E_{z\in B_{\infty}(y,r_{2})}\norm{\nabla f(z)-g(y)}_{1}\leq n^{3/2}\frac{r_{2}}{r_{1}}L
\]
where $g(y)$ is the average of $\nabla f$ over $B_{\infty}(y,r_{2})$.

\begin{figure}[h]

\centering{}\includegraphics[width=0.6\columnwidth]{Figures/Lemma_4\lyxdot 21}
\end{figure}
\end{lem}

\begin{proof}
Let $\omega_{i}(z)=\left\langle \nabla f(z)-g(y),e_{i}\right\rangle $
for all $i\in[n]$. Then, we have that
\[
\int_{B_{\infty}(y,r_{2})}\norm{\nabla f(z)-g(y)}_{1}dz=\sum_{i}\int_{B_{\infty}(y,r_{2})}\left|\omega_{i}(z)\right|dz.
\]
Since $\int_{B_{\infty}(y,r_{2})}\omega_{i}(z)dz=0$, the Poincar\'e
inequality for a box (Theorem \ref{thm:box-poincare} below) shows
that
\begin{align*}
\int_{B_{\infty}(y,r_{2})}\left|\omega_{i}(z)\right|dz & \leq r_{2}\int_{B_{\infty}(y,r_{2})}\norm{\nabla\omega_{i}(z)}_{2}dz\\
 & =r_{2}\int_{B_{\infty}(y,r_{2})}\norm{\nabla^{2}f(z)e_{i}}_{2}dz\\
\sum_{i}\int_{B_{\infty}(y,r_{2})}\left|\omega_{i}(z)\right|dz & \le\sqrt{n}r_{2}\int_{B_{\infty}(y,r_{2})}\norm{\nabla^{2}f(z)}_{F}dz
\end{align*}

Since $f$ is convex, eigenvalues of the Hessian are nonnegative and
so we have 
\[
\norm{\nabla^{2}f(z)}_{F}=\sqrt{\sum_{i}\lambda^{2}_{i}}\le\sum_{i}\lambda_{i}=\tr\nabla^{2}f(z)=\Delta f(z).
\]
Therefore, we have
\begin{align*}
\E_{z\in B_{\infty}(y,r_{2})}\norm{\nabla f(z)-g(y)}_{1} & \leq\sqrt{n}r_{2}\E_{z\in B_{\infty}(y,r_{2})}\Delta f(z)\\
 & =\sqrt{n}r_{2}\Delta h(y)
\end{align*}
where $h=\frac{1}{(2r_{2})^{n}}f\ast\chi_{B_{\infty}(0,r_{2})}$ where
$\chi_{B_{\infty}(0,r_{2})}$ is $1$ on the set $B_{\infty}(0,r_{2})$
and $0$ outside the set.

Integrating by parts, we have that
\[
\int_{B_{\infty}(x,r_{1})}\Delta h(y)dy=\int_{\partial B_{\infty}(x,r_{1})}\left\langle \nabla h(y),n(y)\right\rangle dy
\]
where $\Delta h(y)=\sum_{i}\frac{d^{2}h}{dx^{2}_{i}}(y)$ and $n(y)$
is the outward normal vector on $\partial B_{\infty}(x,r_{1})$ of
the box $B_{\infty}(x,r_{1})$, i.e. a signed standard basis vector.
Since $f$ is $L$-Lipschitz with respect to $\norm{\cdot}_{\infty}$
so is $h$, i.e. $\norm{\nabla h(z)}_{\infty}\leq L$. Hence, we have
that
\[
\E_{y\in B_{\infty}(x,r_{1})}\Delta h(y)\leq\frac{1}{(2r_{1})^{n}}\int_{\partial B_{\infty}(x,r_{1})}\norm{\nabla h(y)}_{\infty}\norm{n(y)}_{1}dy\leq\frac{1}{(2r_{1})^{n}}\cdot2n(2r_{1})^{n-1}\cdot L=\frac{nL}{r_{1}}.
\]
Therefore, we have that
\[
\E_{y\in B_{\infty}(x,r_{1})}\E_{z\in B_{\infty}(y,r_{2})}\norm{\nabla f(z)-g(y)}_{1}\leq n^{3/2}\frac{r_{2}}{r_{1}}L.
\]
\end{proof}
\begin{xca}
For an $L$-gradient-Lipschitz function $f:\R^{n}\rightarrow\R,$
and $h=f*\chi_{B_{\infty}(0,1/2)}$, prove that $h$ is also $L$-gradient-Lipschitz
and $\Delta h(y)=\E_{z\sim B_{\infty}(y,1/2)}(\Delta f(z))$.
\end{xca}

This lemma shows that we can implement an approximate gradient oracle
(GRAD) using an evaluation oracle (EVAL) even for non-differentiable
functions. By the involution property again, this completes all the
reductions in Figure \ref{fig:rel}. With the above fact asserting
that, on average, the gradient is approximated by its average in a
small ball, we now proceed to construct an approximate subgradient,
using only an approximate evaluation oracle. The parameter $r_{2}$
in the algorithm is chosen to optimize the final error of the output.
By making the ratio $r_{2}/r_{1}$ sufficiently small, we can get
a desired error for the subgradient. 

\begin{algorithm2e}[t] 

\caption{$\mathtt{SubgradConvexFunc}(f,x,r_{1},\varepsilon)$}\label{alg:separateFunc}

\SetAlgoLined

\textbf{Require:} $r_{1}>0$, $\norm{\partial f(z)}_{\infty}\leq L$
for any $z\in B_{\infty}(x,2r_{1})$.

Set $r_{2}=\frac{1}{2}\sqrt{\frac{\varepsilon r_{1}}{\sqrt{n}L}}$.

Sample $y\in B_{\infty}(x,r_{1})$ and $z\in B_{\infty}(y,r_{2})$
independently and uniformly at random.

\For{$i=1,2,\cdots,n$ }{

Let $\alpha_{i}$ and $\beta_{i}$ denote the end points of the interval
$B_{\infty}(y,r_{2})\cap\{z+se_{i}:s\in\R\}$.

Set $\tilde{g}_{i}=\frac{f(\beta_{i})-f(\alpha_{i})}{2r_{2}}$ where
we compute $f$ to within $\varepsilon$ additive error.

}

\textbf{Output} $\tilde{g}$ as the approximate subgradient of $f$
at $x$.

\end{algorithm2e}
\begin{lem}
\label{lem:separate_conv_func} Let $r_{1}>0$ and $f$ be a convex
function. Suppose that $\norm{\nabla f(z)}_{\infty}\leq L$ for any
$z\in B_{\infty}(x,2r_{1})$ and suppose that we can evaluate the
function to within $\varepsilon$ additive error for $\varepsilon\leq r_{1}\sqrt{n}L$.
Let $\tilde{g}=\mathtt{SubgradConvexFunc}(f,x,r_{1},\varepsilon)$.
Then, there is a random variable $\zeta\geq0$ with $\E\zeta\leq2\sqrt{\frac{L\varepsilon}{r_{1}}}n^{5/4}$
such that for any $y$
\[
f(y)\geq f(x)+\left\langle \tilde{g},y-x\right\rangle -\zeta\norm{y-x}_{\infty}-4nr_{1}L.
\]
\end{lem}

\begin{proof}
We assume that $f$ is twice differentiable. For general $f$, we
can reduce to this case by viewing it as a limit of twice-differentiable
functions.

\begin{figure}[h]

\begin{centering}
\includegraphics[width=0.6\columnwidth]{Figures/Lemma_4\lyxdot 23}
\par\end{centering}
\end{figure}

First, we assume that we can compute $f$ exactly, namely $\varepsilon=0$.
Fix $i\in[n]$. Let $g(y)$ be the average of $\nabla f$ over $B_{\infty}(y,r_{2})$.
Then, for the function $\tilde{g}$ computed by the algorithm, we
have that
\begin{align*}
\E_{z}\left|\tilde{g}_{i}-g(y)_{i}\right| & =\E_{z}\left|\frac{f(\beta_{i})-f(\alpha_{i})}{2r_{2}}-g(y)_{i}\right|\\
 & \leq\E_{z}\frac{1}{2r_{2}}\int\left|\frac{df}{dx_{i}}(z+se_{i})-g(y)_{i}\right|ds\\
 & =\E_{z}\left|\frac{df}{dx_{i}}(z)-g(y)_{i}\right|
\end{align*}
where we used that both $z+se_{i}$ and $z$ are uniformly distributed
on $B_{\infty}(y,r_{2})$ in the last line. Hence, we have
\[
\E_{z}\norm{\tilde{g}-\nabla f(z)}_{1}\leq\E_{z}\norm{\nabla f(z)-g(y)}_{1}+\E_{z}\norm{\tilde{g}-g(y)}_{1}\leq2\E_{z}\norm{\nabla f(z)-g(y)}_{1}.
\]
Now, applying the convexity of $f$ yields that
\begin{align*}
f(q) & \geq f(z)+\left\langle \nabla f(z),q-z\right\rangle \\
 & =f(z)+\left\langle \tilde{g},q-x\right\rangle +\left\langle \nabla f(z)-\tilde{g},q-x\right\rangle +\left\langle \nabla f(z),x-z\right\rangle \\
 & \geq f(z)+\left\langle \tilde{g},q-x\right\rangle -\norm{\nabla f(z)-\tilde{g}}_{1}\norm{q-x}_{\infty}-\norm{\nabla f(z)}_{\infty}\norm{x-z}_{1}.
\end{align*}
Now, using $f$ is $L$-Lipschitz between $x$ and $z$, we have that
$f(z)\geq f(x)-L\cdot\norm{x-z}_{1}$. Hence, we have
\[
f(q)\geq f(x)+\left\langle \tilde{g},q-x\right\rangle -\norm{\nabla f(z)-\tilde{g}}_{1}\norm{q-x}_{\infty}-2L\norm{x-z}_{1}.
\]
Note that $\norm{x-z}_{1}\leq n\cdot\norm{x-z}_{\infty}\leq n(r_{1}+r_{2})$
by assumption. Moreover, we can apply Lemma~\ref{lem:convex_almost_flat}
to bound $\norm{\nabla f(z)-\tilde{g}}_{1}$ and use $r_{2}=\frac{1}{2}\sqrt{\frac{\varepsilon r_{1}}{\sqrt{n}L}}\leq r_{1}$
to get
\[
f(q)\geq f(x)+\left\langle \tilde{g},q-x\right\rangle -\zeta\norm{q-x}_{\infty}-4nr_{1}L
\]
with $\E\zeta\leq2n^{3/2}\frac{r_{2}}{r_{1}}L$.

Since we only compute $f$ up to $\varepsilon$ additive error, it
introduces $\frac{\varepsilon}{r_{2}}$ additive error in $\tilde{g}_{i}$.
Hence, we instead have that
\[
\E\zeta\leq2n^{3/2}\frac{r_{2}}{r_{1}}L+\frac{\varepsilon n}{2r_{2}}.
\]
Setting $r_{2}=\frac{1}{2}\sqrt{\frac{\varepsilon r_{1}}{\sqrt{n}L}}$
completes the proof.
\end{proof}
Note that if we happened to have an exact oracle for $f$, then we
can make $r_{2}$ arbitrarily small.
\begin{xca}
Suppose in Algorithm \ref{alg:separateFunc} we apply the idea of
a linear approximation directly to a small box around the query point
$x$, i.e., we pick a random point $z$, avoiding the intermediate
point $y$ entirely, and use $z$ to find endpoints of chords through
$z$ and compute an estimate of the gradient. Why is this not accurate
in general?
\end{xca}

\begin{xca}
What is the best possible bound in Lemma \ref{lem:convex_almost_flat},
in particular, is the dependence on the dimension tight?
\end{xca}

In the proof of Lemma \ref{lem:convex_almost_flat}, we used the following
fact applied to the case when $\Omega$ is a cube. In this case, the
coefficient on the RHS is given by the Cheeger or KLS constant of
the cube and is $1$. This is an example of an isoperimetric inequality.
Such inequalities will play an important role later in this book.
\begin{thm}[$L^{1}$-Poincar\'e inequality]
\label{thm:box-poincare} Let $\Omega$ be connected, bounded and
open. Then the following (best-possible) inequality holds for any
smooth function $f:\Omega\rightarrow\R$:
\[
\norm{f-\frac{1}{\left|\Omega\right|}\int_{\Omega}f(x)\,dx}_{L^{1}(\Omega)}\le\left(\sup_{S\subset\Omega}\frac{2|S||\Omega\setminus S|}{|\partial S||\Omega|}\right)\norm{\nabla f}_{L^{1}(\Omega)}
\]
 where the supremum is over all subsets $S$ such that $S$ and $\Omega\setminus S$
are both connected.
\end{thm}

\begin{xca}
Prove the inequality in Theorem \ref{thm:box-poincare} using the
classical co-area formula.
\end{xca}


\section{Separation via Membership \label{sec:From-Membership-to}}

In this section, we show how to implement a separation oracle for
a convex set using a nearly linear number of queries to a membership
oracle. We divide this into two steps. In the first step (Algorithm~\ref{alg:separateFunc}),
we compute an approximate subgradient of a given Lipschitz convex
function. Using this, in the second step (Algorithm~\ref{alg:sep_set}),
we compute an approximate separating hyperplane. 
\begin{thm}
Let $K$ be a convex body such that $B(0,1)\subseteq K\subseteq B(0,R)$.
Then, for any $0<\eta\leq1/2$, we can compute an $\eta$-approximate
separation oracle for K using $O(n\log(nR/\eta))$ queries to a membership
oracle for $K$.
\end{thm}

Throughout this section, let $K\subseteq\R^{n}$ be a convex set that
contains $B_{2}(0,r)$ and is contained in $B_{2}(0,R)$. Recall that
$B_{p}(x,r)$ is the $p$-norm ball of radius $r$ centered at $x$.
By an approximate separation oracle we mean that either the queried
point $x$ lies within distance $\eta$ of $K$, or the oracle provides
a halfspace $c^{T}y\le c^{T}x+\eta$ for all points $y\in K$ at distance
at least $\eta$ from the boundary of $K$. We also define 
\[
B_{p}(K,-\delta)\defeq\{x\in\R^{n}:B_{p}(x,\delta)\subseteq K\}.
\]
Given some point $x\notin K$, we wish to separate $x$ from $K$
using a halfspace. To do this, we reduce this problem to computing
an approximate subgradient of a Lipschitz convex function $h_{x}(d)$
defined for points in $K.$ Roughly speaking, it is the ``height''
(or distance from the boundary) of a point $d$ in the direction of
$x$. Let 
\[
\alpha_{x}(d)=\max_{d+\alpha x\in K}\alpha\mbox{ \quad and \quad}h_{x}(d)=-\alpha_{x}(d)\norm x_{2}.
\]
 Note that $d+\alpha_{x}(d)x$ is the last point in $K$ on the line
through $d\in K$ in the direction of $x$, and $-h_{x}(d)$ is the
$\ell_{2}$ distance from this boundary point to $d$ (see Fig.\ref{fig:height-function}).

\begin{figure}[h]
\begin{centering}
\includegraphics[width=0.6\columnwidth]{Figures/Thrm_4\lyxdot 27}
\par\end{centering}
\centering{}\caption{The convex height function $h_{x}$}
\label{fig:height-function}
\end{figure}

\begin{algorithm2e}[h]

\caption{$\mathtt{Separate}_{\varepsilon,\rho}(K,x)$}\label{alg:sep_set}

\SetAlgoLined

\textbf{Require:} $B_{2}(0,r)\subset K\subset B_{2}(0,R)$.

\uIf{$\text{MEM}_{\varepsilon}(K)$ asserts that $x\in B(K,\epsilon)$}{

\textbf{Output:} ``$x\in B(K,\varepsilon)$''.

}\ElseIf{$x\notin B_{2}(0,R)$}{

\textbf{Output:} the halfspace $\{y:0\geq\left\langle y-x,x\right\rangle \}$.

}

Let $\kappa=R/r$, $\alpha_{x}(d)=\max_{d+\alpha x\in K}\alpha$ and
$h_{x}(d)=-\alpha_{x}(d)\norm x_{2}$. 

The evaluation oracle of $\alpha_{x}(d)$ can be implemented via binary
search and $\text{MEM}_{\varepsilon}(K)$. 

Compute $\tilde{g}=\mathtt{SubgradConvexFunc}(h_{x},0,r_{1},4\varepsilon)$
with $r_{1}=n^{1/6}\varepsilon^{1/3}R^{2/3}\kappa^{-1}$ and the evaluation
oracle of $\alpha_{x}(d)$.

\textbf{Output:} the halfspace
\[
\left\{ y:\frac{31}{\rho}n^{7/6}R^{2/3}\kappa\varepsilon^{1/3}\geq\left\langle \tilde{g},y-x\right\rangle \right\} 
\]

\end{algorithm2e}

The output of the algorithm for separation is a halfspace that approximately
contains $K$, and the input point $x$ is close to its bounding hyperplane.
It uses a call to the subgradient function above. 

We now proceed to analyze the height function.
\begin{lem}
\label{lem:h_convex} $h_{x}(d)$ is convex on $K$.
\end{lem}

\begin{proof}
Let $d_{1},d_{2}\in K$ and $\lambda\in[0,1]$. Now $d_{1}+\alpha_{x}(d_{1})x\in K$
and $d_{2}+\alpha_{x}(d_{2})x\in K$. Consequently,
\[
\left[\lambda d_{1}+(1-\lambda)d_{2}\right]+\left[\lambda\cdot\alpha_{x}(d_{1})+(1-\lambda)\cdot\alpha_{x}(d_{2})\right]x\in K\,.
\]
Therefore, if we let $d\defeq\lambda d_{1}+(1-\lambda)d_{2}$ we see
that $\alpha_{x}(d)\geq\lambda\cdot\alpha_{x}(d_{1})+(1-\lambda)\cdot\alpha_{x}(d_{2})$
and $h_{x}(\lambda d_{1}+(1-\lambda)d_{2})\leq\lambda h_{x}(d_{1})+(1-\lambda)h_{x}(d_{2})$
as claimed.
\end{proof}
\begin{lem}
\label{lem:h_Lip} $h_{x}$ is $\left(\frac{R+\delta}{r-\delta}\right)$-Lipschitz
over points in $B_{2}(0,\delta)$ for any $\delta<r$. 

\begin{figure}[h]

\begin{centering}
\includegraphics[width=0.4\columnwidth]{Figures/Lemma_4\lyxdot 29}
\par\end{centering}
\end{figure}
\end{lem}

\begin{proof}
Let $d_{1},d_{2}$ be arbitrary points in $B(0,\delta)$. We wish
to upper bound $\left|h_{x}(d_{1})-h_{x}(d_{2})\right|$ in terms
of $\norm{d_{1}-d_{2}}_{2}$. We assume without loss of generality
that $\alpha_{x}(d_{1})\geq\alpha_{x}(d_{2})$ and therefore
\[
\left|h_{x}(d_{1})-h_{x}(d_{2})\right|=\left|\alpha_{x}(d_{1})\norm x_{2}-\alpha_{x}(d_{2})\norm x_{2}\right|=\left(\alpha_{x}(d_{1})-\alpha_{x}(d_{2})\right)\norm x_{2}\,.
\]
Consequently, it suffices to lower bound $\alpha_{x}(d_{2})$. We
split the analysis into two cases.

Case 1: $\norm{d_{2}-d_{1}}_{2}\geq r-\delta$. Since $0\geq h_{x}(d_{1}),h_{x}(d_{2})\geq-R-\delta$,
we have that 
\[
\left|h_{x}(d_{1})-h_{x}(d_{2})\right|\leq R+\delta\leq\frac{R+\delta}{r-\delta}\norm{d_{2}-d_{1}}_{2}.
\]

Case 2: $\norm{d_{2}-d_{1}}_{2}\leq r-\delta$. If $d_{1}=d_{2}$,
the claim is trivial. Otherwise, we consider the point $d_{3}=d_{1}+\frac{d_{2}-d_{1}}{\lambda}$
with $\lambda=\norm{d_{2}-d_{1}}_{2}/(r-\delta)$. Note that 
\[
\norm{d_{3}}_{2}\leq\norm{d_{1}}_{2}+\frac{1}{\lambda}\norm{d_{2}-d_{1}}_{2}\leq\delta+\frac{1}{\lambda}\norm{d_{2}-d_{1}}_{2}\leq r.
\]
Hence, $d_{3}\in K$. Since $\lambda\in[0,1]$ and $K$ is convex,
we have that $\lambda\cdot d_{3}+(1-\lambda)\cdot\left[d_{1}+\alpha_{x}(d_{1})x\right]\in K$.
Now, we note that
\[
\lambda\cdot d_{3}+(1-\lambda)\cdot\left[d_{1}+\alpha_{x}(d_{1})x\right]=d_{2}+(1-\lambda)\cdot\alpha_{x}(d_{1})x
\]
and this shows that 
\[
\alpha_{x}(d_{2})\geq(1-\lambda)\cdot\alpha_{x}(d_{1})=\left(1-\frac{\norm{d_{2}-d_{1}}_{2}}{r-\delta}\right)\cdot\alpha_{x}(d_{1}).
\]
Since $d_{1}+\alpha_{x}(d_{1})x\in K\subset B_{2}(0,R)$, we have
that $\alpha_{x}(d_{1})\cdot\norm x_{2}\leq R+\delta$ and hence
\[
\left|h_{x}(d_{1})-h_{x}(d_{2})\right|=(\alpha_{x}(d_{1})-\alpha_{x}(d_{2}))\cdot\norm x_{2}\leq\alpha_{x}(d_{1})\cdot\norm x_{2}\frac{\norm{d_{2}-d_{1}}_{2}}{r-\delta}\leq\frac{R+\delta}{r-\delta}\norm{d_{2}-d_{1}}_{2}.
\]
 In either case, as claimed we have
\[
\left|h_{x}(d_{1})-h_{x}(d_{2})\right|\leq\frac{R+\delta}{r-\delta}\norm{d_{2}-d_{1}}_{2}.
\]
\end{proof}
The next lemma shows that $h_{x}$ gives us a way to implement an
approximate separation oracle and only needs access to an approximate
evaluation oracle for $h_{x}$ which in turn only needs an approximate
membership oracle for $K$. To be precise, we define approximate oracles.
\begin{defn}[Separation Oracle (SEP)]
\label{def:SEP-1} Queried with a vector $y\in\Rn$ and real numbers
$\delta,\delta'>0$, with probability at least $1-\delta'$, the oracle
either

\begin{itemize}
\item assert that $y\in B(K,\delta)$, or
\item find a unit vector $c\in\Rn$ such that $c^{T}x\leq c^{T}y+\delta$
for all $x\in B(K,-\delta)$.
\end{itemize}
We let $\text{SEP}_{\delta,\delta'}(K)$ be the time complexity of
this oracle.
\end{defn}

\begin{defn}[Membership Oracle (MEM)]
\label{def:MEM-1} Queried with a vector $y\in\Rn$ and real numbers
$\delta,\delta'>0$, with probability at least $1-\delta'$, either

\begin{itemize}
\item assert that $y\in B(K,\delta)$, or
\item assert that $y\notin B(K,-\delta)$.
\end{itemize}
We let $\text{MEM}_{\delta,\delta'}(K)$ be the time complexity of
this oracle.
\end{defn}

We can now state the main lemma for the separation oracle. Since the
algorithm is randomized, we have a parameter $\rho\in(0,1)$ to denote
the probability of failure.
\begin{lem}
\label{lem:separate_set} Let $K$ be a convex set satisfying $B_{2}(0,r)\subset K\subset B_{2}(0,R)$.
Let $0<\rho<1$ and $\varepsilon>0$ be sufficiently small so that
$r_{1}\le r/4$. Then, with probability $1-\rho$, $\mathtt{Separate}_{\varepsilon,\rho}(K,x)$
outputs a halfspace that contains $K$. 
\end{lem}

\begin{proof}
When $x\notin B_{2}(0,R)$, the algorithm outputs a valid separation
for $B_{2}(0,R)$. For the rest of the proof, we assume $x\notin B(K,-\varepsilon)$
(due to the membership oracle) and $x\in B_{2}(0,R)$.

By Lemma \ref{lem:h_convex} and Lemma \ref{lem:h_Lip}, $h_{x}$
is convex with Lipschitz constant $3\kappa$ on $B_{2}(0,\frac{r}{2})$.
By our assumption on $\varepsilon$ and our choice of $r_{1}$, we
have that $B_{\infty}(0,2r_{1})\subset B_{2}(0,\frac{r}{2})$. Hence,
we can apply Lemma \ref{lem:separate_conv_func} to get that
\begin{align}
h_{x}(y) & \geq h_{x}(0)+\left\langle \tilde{g},y\right\rangle -\zeta\norm y_{\infty}-12nr_{1}\kappa\label{eq:hxy-1}
\end{align}
for any $y\in K$. Note that $-\frac{x}{\kappa}\in K$ and $h_{x}(-\frac{x}{\kappa})=h_{x}(0)-\frac{1}{\kappa}\norm x_{2}$.
Hence, we have
\[
h_{x}(0)-\frac{1}{\kappa}\norm x_{2}=h_{x}(-\frac{1}{\kappa}x)\geq h_{x}(0)+\left\langle \tilde{g},-\frac{1}{\kappa}x\right\rangle -\frac{1}{\kappa}\zeta\norm x_{\infty}-12nr_{1}\kappa.
\]
Therefore, we have
\begin{equation}
\left\langle \tilde{g},x\right\rangle \geq\norm x_{2}-\zeta\norm x_{\infty}-12nr_{1}\kappa^{2}.\label{eq:lowerbound_g}
\end{equation}
Now, we note that $x\notin B(K,-\varepsilon)$. Using that $B(0,r)\subset K$,
we have $(1-\frac{\varepsilon}{r})K\subset B(K,-\varepsilon)$. Hence,
\[
h_{x}(0)\geq-\frac{\norm x_{2}}{\left(1-\frac{\varepsilon}{r}\right)}\geq-\norm x_{2}-\varepsilon\kappa.
\]
Therefore, we have
\[
h_{x}(0)+\left\langle \tilde{g},x\right\rangle \geq-\zeta\norm x_{\infty}-12nr_{1}\kappa^{2}-\varepsilon\kappa
\]
Combining this with (\ref{eq:hxy-1}), we have that 
\begin{align*}
h_{x}(y) & \geq\left\langle \tilde{g},y-x\right\rangle -\zeta\norm y_{\infty}-\zeta\norm x_{\infty}-12nr_{1}\kappa-12nr_{1}\kappa^{2}-\varepsilon\kappa\\
 & \geq\left\langle \tilde{g},y-x\right\rangle -2\zeta R-24nr_{1}\kappa^{2}-\varepsilon\kappa
\end{align*}
for any $y\in K$. Recall from Lemma \ref{lem:separate_conv_func}
that $\zeta$ is a positive random scalar independent of $y$ satisfying
$\E\zeta\leq2\sqrt{\frac{3\kappa\varepsilon}{r_{1}}}n^{5/4}$. For
any $y\in K$, we have that $h_{x}(y)\leq0$ and hence $\tilde{\zeta}\geq\left\langle \tilde{g},y-x\right\rangle $
where $\tilde{\zeta}$ is a random scalar independent of $y$ satisfying
\begin{align*}
\E\tilde{\zeta} & \le4\sqrt{\frac{3\kappa\varepsilon}{r_{1}}}n^{5/4}R+24nr_{1}\kappa^{2}+\varepsilon\kappa\\
 & \leq31n^{7/6}R^{2/3}\varepsilon^{1/3}\kappa.
\end{align*}
where we used $r_{1}=n^{1/6}\varepsilon^{1/3}R^{2/3}/\kappa$ and
$0\leq\varepsilon\leq r$. The result then follows using Markov's
inequality. 
\end{proof}
\begin{xca}
Suppose we can evaluate the subgradient of $h_{x}$ exactly for a
convex set $K$ containing the origin. Give a short proof that for
any $x\not\in K$, we have $\langle\nabla h_{x}(0),y-x\rangle\le0$
for all $y\in K$. 
\end{xca}

\begin{thm}
\label{thm:separate_set} Let $K$ be a convex set satisfying $B_{2}(0,1/\kappa)\subset K\subset B_{2}(0,1)$.
For any $0<\eta<\frac{1}{2}$, we have that
\[
\text{SEP}_{\eta}(K)\leq O\left(n\log\left(\frac{n\kappa}{\eta}\right)\right)\text{MEM}_{(\eta/n\kappa)^{O(1)}}(K).
\]
\end{thm}

\begin{proof}
First, we bound the running time. Note that the bottleneck is to compute
$h_{x}$ with $\delta$ additive error. Since $-O(1)\leq h_{x}(y)\leq0$
for all $y\in B_{2}(0,O(1))$, one can compute $h_{x}(y)$ by binary
search with $O(\log(1/\delta))$ calls to the membership oracle.

Next, we check that $\mathtt{Separate}_{\delta,\rho}(K,x)$ is indeed
a separation oracle. Note that $\tilde{g}$ may not be a unit vector
and we need to re-normalize $\tilde{g}$ by $1/\norm{\tilde{g}}_{2}$.
So, we need a lower bound $\norm{\tilde{g}}_{2}$. 

From (\ref{eq:lowerbound_g}) and our choice of $r_{1}$, if we set
$\varepsilon=\delta$, where $\delta$ is the error and $\rho$ is
the failure probability of the subgradient oracle, then we have that
\begin{align*}
\left\langle \tilde{g},x\right\rangle  & \geq\norm x_{2}-\zeta\norm x_{\infty}-12nr_{1}\kappa^{2}\geq\frac{r}{4}.
\end{align*}
Hence, we have that $\norm{\tilde{g}}_{2}\geq\frac{1}{4\kappa}$.
Therefore, this algorithm is a separation oracle with error $\frac{400}{\rho}n^{7/6}\kappa^{2}\delta^{1/3}$
and by the union bound over all calls to the membership oracle, with
failure probability $O(\rho+n\log(1/\delta)\delta)$. Therefore,
\[
\text{SEP}_{O(\max\{n^{7/6}\kappa^{2}\delta^{1/3}/\rho,\rho,n\log(1/\delta)\delta\})}(K)\leq O(\log(1/\delta))\text{MEM}_{\delta}(K).
\]
Setting $\rho=\sqrt{n^{7/6}\kappa^{2}\delta^{1/3}}$ and $\delta=\Theta\left(\frac{\eta^{6}}{n^{7/2}\kappa^{6}}\right)$,
we have that 
\[
\text{SEP}_{\eta}(K)\leq O(\log(\frac{n\kappa}{\eta}))\text{MEM}_{\eta^{6}/(n^{7/2}\kappa^{6})}(K).
\]
\end{proof}

\subsection{Gradient from Evaluation via Auto Differentiation }

In this section, we give yet another way to compute gradients using
evaluations.
\begin{thm}
\label{thm:auto_diff} Given a function $f:\R^{n}\rightarrow\R$ represented
by a circuit whose gates compute differentiable functions of a finite
number of variables. Suppose that $f(x)$ can be computed in time
$\mathcal{T}$ (namely, the circuit has $\mathcal{T}$ edges). Then
we can compute $\nabla f(x)$ in $O(\mathcal{T})$ time.
\end{thm}

\begin{rem*}
In practice, the runtime is roughly $2\mathcal{T}$ assuming we have
enough memory. Check out JAX for a modern implementation.

Before proving it formally, we first go through an example. Consider
the function $f(x_{1},x_{2})=\sin(x_{1}/x_{2})+x_{1}x_{2}$. We use
$x_{i}$ to denote both the input and all intermediate variables.
Then, we can write the program in $\mathcal{T}=6$ steps:
\end{rem*}
\begin{itemize}
\item $x_{1}=x_{1}$, $x_{2}=x_{2}$, $x_{3}=x_{1}/x_{2}$, $x_{4}=\sin(x_{3})$,
$x_{5}=x_{1}x_{2}$, Output $x_{6}=x_{4}+x_{5}$.
\end{itemize}
Note that each step involves computing
\[
x_{i}=f_{i}(\overset{\text{only a few }x_{j}}{\overbrace{x_{1},\cdots,x_{i-1}}})
\]
with simple functions $f_{i}$ whose derivatives we know how to compute.
The key idea is to compute $\frac{\partial f}{\partial x_{i}}$ not
just for the inputs $x_{1}$ and $x_{2}$, but also for all intermediate
variables. Here, we use $\frac{\partial f}{\partial x_{i}}$ to denote
the derivative of $f$ with respect to $x_{i}$ while fixing $x_{1},x_{2},\cdots,x_{i-1}$
(and other inputs if $x_{i}$ is an input). For the example above,
suppose we want to compute $\nabla f(\pi,2)$, we can first compute
all $x_{i}$ from $i=1,2,\cdots,6$, then $\frac{\partial f}{\partial x_{i}}$
in the reverse order from $i=6,5,\cdots,1$:
\begin{itemize}
\item $x_{1}=\pi$, $x_{2}=2$, $x_{3}=\pi/2$, $x_{4}=\sin(x_{3})=1$,
$x_{5}=x_{1}x_{2}=2\pi$, $x_{6}=x_{4}+x_{5}=2\pi+1$.
\item $\frac{\partial f}{\partial x_{6}}=1$, $\frac{\partial f}{\partial x_{5}}=\frac{\partial\left(x_{4}+x_{5}\right)}{\partial x_{5}}=1$,
$\frac{\partial f}{\partial x_{4}}=\frac{\partial\left(x_{4}+x_{5}\right)}{\partial x_{4}}=1$,
\item $\frac{\partial f}{\partial x_{3}}=\frac{\partial f}{\partial x_{4}}\frac{\partial x_{4}}{\partial x_{3}}=1\cdot\cos(x_{3})=0$,
\item $\frac{\partial f}{\partial x_{2}}=\frac{\partial f}{\partial x_{3}}\frac{\partial x_{3}}{\partial x_{2}}+\frac{\partial f}{\partial x_{5}}\frac{\partial x_{5}}{\partial x_{2}}=0\cdot(-\frac{x_{1}}{x^{2}_{2}})+1\cdot x_{1}=\pi$,
\item $\frac{\partial f}{\partial x_{1}}=\frac{\partial f}{\partial x_{3}}\frac{\partial x_{3}}{\partial x_{1}}+\frac{\partial f}{\partial x_{5}}\frac{\partial x_{5}}{\partial x_{1}}=0\cdot(\frac{1}{x_{2}})+1\cdot x_{2}=2$.
\end{itemize}
The general case is similar. See $\mathtt{AutoDifferentiation}$ for
the algorithm.

\begin{algorithm2e}[H]

\caption{$\mathtt{AutoDifferentiation}$}

\SetAlgoLined

\textbf{Input:} a function $f(x_{1},x_{2},\cdots,x_{n})$ given by
$f(x_{1},x_{2},\cdots,x_{n})=x_{m}$ and
\[
x_{i}=f_{i}(x_{1},\cdots,x_{i-1})\text{ for }i=n+1,n+2,\cdots,m
\]

\For{$i=n+1,n+2,\cdots,m$}{

Compute $x_{i}=f_{i}(x_{1},\cdots,x_{i-1})$.

}

Let $\frac{\partial f}{\partial x_{m}}=1$.

\For{$i=m-1,\cdots,1$}{

Let $L_{i}$ be the set of $j$ such that $f_{j}$ depends on $x_{i}$
(i.e. $x_{j}$ directly depends on $x_{i}$).

Compute $\frac{\partial f}{\partial x_{i}}=\sum_{j\in L_{i}}\frac{\partial f}{\partial x_{j}}\frac{\partial x_{j}}{\partial x_{i}}$.

}

\end{algorithm2e}
\begin{proof}[Proof of Theorem \ref{thm:auto_diff}]
We prove by induction that the formula $\frac{\partial f}{\partial x_{i}}=\sum_{j\in L_{i}}\frac{\partial f}{\partial x_{j}}\frac{\partial x_{j}}{\partial x_{i}}$
is correct. For the base case $i=m$, we have $f=x_{m}$ and hence
$\frac{\partial f}{\partial x_{m}}=1$. For the induction, we let
$L_{i}=\{j_{1},j_{2},\cdots,j_{k}\}$. If we fix variables $x_{1},x_{2},\cdots,x_{i-1}$,
then $f$ is a function of $x_{i}$ (and of other inputs if $x_{i}$
is an input). Since only $x_{j_{1}},x_{j_{2}},\cdots,x_{j_{k}}$ depend
on $x_{i}$, we can also view $f$ as a function of $x_{j_{1}},x_{j_{2}},\cdots,x_{j_{k}}$.
Applying the chain rule to these direct dependencies, we have
\[
\frac{\partial f}{\partial x_{i}}=\sum_{j\in L_{i}}\frac{\partial f}{\partial x_{j}}\frac{\partial x_{j}}{\partial x_{i}}.
\]

To bound the runtime, we define the computation graph $G$ to be a
graph on $x_{1},x_{2},\cdots,x_{m}$ such that $i\rightarrow j$ if
$f_{j}$ depends on $x_{i}$. Note that each edge is examined $O(1)$
times whether evaluating $f$ or its gradient. Hence, the cost of
computing $f$ and the cost of our algorithm are both $\Theta(\mathcal{T})$
where $\mathcal{T}$ is the number of edges in $G$. This completes
the proof. 
\end{proof}
We note that an efficient implementation of the chain rule is the
heart of the backpropagation algorithm for neural networks. To conclude
this section, we see that Theorem \ref{thm:auto_diff} can be surprisingly
useful even for some simple explicit functions.
\begin{cor}
If we can compute $f(A)=\det A$ exactly in time $\mathcal{T}$ for
nonsingular $A$, then we can compute $A^{-1}$ exactly in $O(\mathcal{T})$.
\end{cor}

\begin{proof}
Note that $\frac{\partial}{\partial A_{ij}}\det A=\text{adj}(A)_{ji}=\det A\cdot(A^{-1})_{ji}$.
Hence, $\nabla\log\det A=A^{-\top}$. Theorem \ref{thm:auto_diff}
shows that computing $A^{-\top}$ can be done as fast as $\det A$.
\end{proof}

\subsection{Gradient from Evaluation via Complex Step Differentiation}

Auto differentiation is great if the function can be computed exactly.
However, it does not work well if $f$ can only be approximated or
the computation of $f$ involves too many variables. For example,
if $f$ is the energy usage of some aircraft design, then $f$ can
only be computed approximately via simulation and such computation
is memory expensive (if we need to store all intermediate variables)
and not exact, so the method in the previous subsection is not practical.

To introduce complex step differentiation, we recall the formula of
finite difference:
\begin{itemize}
\item Forward/Backward difference: $\frac{f(x\pm h)-f(x)}{\pm h}=f'(x)+hf''(\zeta)$.
\item Central difference: $\frac{f(x+h)-f(x-h)}{2h}=f'(x)+O(h^{2})f'''(\zeta)$.
\end{itemize}
Note that the error analysis above assumes no numerical error involved.
The formula involves subtracting two close numbers and dividing by
a small number and this step creates lots of error. If we are computing
$f$ as floating point, we have
\[
\frac{(1\pm\epsilon)f(x\pm h)-(1\pm\epsilon)f(x)}{\pm h}=f'(x)+O(\frac{\epsilon}{h})f(\zeta_{1})+O(h)f''(\zeta_{2})
\]
\[
\frac{(1\pm\epsilon)f(x+h)-(1\pm\epsilon)f(x-h)}{2h}=f'(x)+O(\frac{\epsilon}{h})f(\zeta_{1})+O(h^{2})f'''(\zeta_{2})
\]
where $\epsilon$ is the floating point precision. Suppose that $f,f',f'',f'''$
are all bounded by constants and suppose $\epsilon=(10)^{-8}$ (single
precision floating point), then we should pick $h=\sqrt{\epsilon}$
for the first case and $h=\epsilon^{1/3}$ in the second case. Hence,
forward/backward difference gives only accuracy $\epsilon^{1/2}\approx(10)^{-4}$
and the central difference gives only accuracy $\epsilon^{2/3}\approx(10)^{-5}$.
In reality, the error will be larger because $f''$ and $f'''$ usually
are large.

The complex step differentiation uses the step
\[
\frac{\mathrm{Im}f(x+ih)}{h}\approx f'(x)+O(h^{2})f'''(\zeta).
\]
This formula works only for complex analytic functions, but it avoids
subtracting numbers. In practice, this really allows us to compute
$f'(x)$ close to machine accuracy. In general, ensuring algorithms
give close to machine accuracy is important because algorithms often
stack on top of each other.
